\section{General Conclusion}


Improving constraints on the growth rate of large-scale structures, 
$f\sigma_8$, is one of the central challenges of modern cosmology, as it 
provides a direct test of gravity on cosmological scales. During this 
PhD, I addressed this challenge along two complementary directions using 
Type Ia supernovae from the ZTF survey: improving the precision and 
characterization of supernova distance measurements, and developing a 
field-level inference approach to reconstruct the density and velocity 
fields of the nearby Universe. \\

The first direction focused on improving supernova distance measurements 
through the continued development and application of the NaCl framework, 
an algorithm designed to train empirical Type Ia supernova 
spectrophotometric models while simultaneously fitting their light 
curves. In its current implementation, the trained model adopts a 
SALT2-like parameterization consisting of two global spectral surfaces, 
$M_0(p,\lambda)$ and $M_1(p,\lambda)$, together with a color law, 
$CL(\lambda)$. 

After its initial development, NaCl required a substantial rewrite to 
make it suitable for \textsc{Lemaître}-scale datasets, and I contributed to this  
effort. Among the developments introduced during this rewrite, I 
extended NaCl to handle spectrophotometric data directly. In particular, 
the use of the SNfactory dataset required adapting the framework so that 
accurately flux-calibrated spectra could be modeled without introducing 
the spectral recalibration polynomial normally used for heterogeneous 
spectroscopic data. This also required reintroducing the color law in 
the spectral model and fitting the $X_0$ parameter of each supernova 
directly from its spectra. Applied to the SNfactory dataset, NaCl 
successfully trained a model on real spectrophotometric data and 
revealed some of the limitations of the SALT2 parameterization in 
describing supernova diversity, in particular in the UV and around the 
Si II and Ca II spectral features.

My work also included the development of an adaptive regularization 
scheme designed to strengthen the model constraints in regions of 
phase-wavelength space where the training data are sparse. This scheme 
is not currently part of the production version of NaCl, as its present 
implementation prevents the minimization from fully converging. Further 
work will therefore be required to stabilize the procedure, possibly 
through a reformulation of the regularization term. 

NaCl is also an essential component of the \textsc{Lemaître} project, which 
combines data from the ZTF, SNLS, and HSC surveys to construct an 
independent Hubble diagram containing approximately 3000 Type Ia 
supernovae over the redshift range $0<z<1.5$. The development of the new 
georges cosmological analysis pipeline within \textsc{Lemaître} required an 
extensive programme of validation on simulated datasets through a series 
of data challenges. The pre-DC1 and DC1 exercises presented in this 
thesis were used to validate the NaCl component and its propagation 
through the downstream analysis. 

These tests show that NaCl is able to reproduce the simulated 
photometric and spectroscopic observations accurately and to recover the 
supernova parameters and their covariance structure with good precision. 
The uncertainties inferred from the Fisher matrix are overall consistent 
with the empirical uncertainties measured across repeated realizations, 
although residual discrepancies remain for some parameters and require 
further investigation. 

However, once the complete NaCl-EDRIS-\texttt{cosmologix} chain is considered, 
significant biases are observed in the reconstructed distance moduli, 
particularly near the transitions between the different survey samples. 
These biases propagate into the inferred cosmological parameters. Their 
origin has not yet been conclusively identified and may involve model 
regularization, effective selection effects introduced by the analysis 
cuts, or limitations in the propagation of the covariance matrix. 
Resolving these effects is an essential step before the pipeline can be 
applied to the blinded \textsc{Lemaître} dataset. \\


The second front focused on applying the BORG Bayesian field-level inference pipeline on ZTF Type Ia supernovae. The goal of this is to obtain the most accurate reconstruction 
of the velocity field from the supernovae by directly inferring the initial conditions of the underlying matter field, and with the use of a gravity model, obtain the 
large scale structures. BORG has proven to be more accurate than other widely used methods in reconstructing the velocity field. With the development of a novel BORG-PV pipeline 
that includes an improved selection modeling and supernova prior, tests on self-consistent simulations have shown the capability of accurate reconstruction of the density and velocity 
fields. The pipeline has also proved capable in inferring cosmological parameters such as $\Omega_m$ and $\sigma_8$ under the assumption of $\Lambda$CDM. 
When applied to more realistic ZTF-like simulations generated from the 
Uchuu N-body simulation, the pipeline has proven to struggle in reconstructing the 
density and velocity fields. Further investigation has shown that may be due to model mismatches. 
These results motivate further investigation of the simulation framework and supernova modeling before applying the pipeline to real ZTF data.
These results pave the way for future analyses for the measurement of $f\sigma_8$ and for cross-survey analyses with galaxy clustering. 



\section{Future Prospects}

The work conducted during this PhD presents the base of a larger and more complete analysis of the ZTF Type Ia supernova data. On the side of NaCl, 
additional testing on capturing calibration uncertainties is needed within DC3 and a color scatter implementation is needed for DC4 before finalizing the NaCl 
framework.  
Moreover, a more complex 
error model needs to be implemented in order to capture uncertainties that are wavelength-phase-dependent as well as to spot outliers that may have 
contaminated the training dataset. 
Once that is complete, DC5 can be tested and the pipeline can operate on the full-blinded \textsc{Lemaître} dataset 
in order to obtain 
an independent measurement of the dark energy equation of state $w$ 
and its potential time variation $w_a$, with a dataset of $\sim 3000$ 
Type Ia supernovae spanning $0 < z < 1.5$. 
Running the full pipeline will also give us access to accurate measurement of the supernovae distances, which will then allow us to proceed with $f\sigma_8$ measurements. \\


On the BORG-PV side, more realistic ZTF simulations are needed in order to validate the pipeline. This includes applying masks to only cover the ZTF footprint 
as well as masking the Milky Way on the line of sight. 
An important aspect not yet accounted for is the correlation between 
supernova distance measurements, which were seen in Chapter \ref{chap4}. 
Accounting for these correlations within the BORG-PV likelihood is a challenge that must be met before operating on the ZTF data. 
Lastly, in order to operate on the full ZTF dataset, a more complex selection function is needed within the likelihood. 

Another aspect that can be explored within BORG-PV is the possibility of conducting a cross-survey analysis. 
One approach would be to use the Manticore-Local constrained realization 
of the local Universe as the density field for BORG-PV. 
This would enable a joint inference 
combining the velocity information from ZTF supernovae with the density 
field information from the galaxy survey, effectively performing a 
multi-tracer field-level inference. 
The exact incorporation of a multi-probe analysis is yet to be determined, and extensive testing on simulations is required before moving on to applications on real data. \\

The methods developed during this PhD lay the groundwork for future 
analyses with the Vera C. Rubin Observatory's LSST, which will provide 
an order of magnitude more Type Ia supernovae at low redshift, enabling 
tighter constraints on $f\sigma_8$ and a rigorous test of $\Lambda$CDM.



